% Clase
\documentclass[11pt,a4paper,spanish,twoside]{book}

% Órdenes auxiliares
% Español
\usepackage[spanish]{babel}
\usepackage[utf8]{inputenc}
\usepackage[T1]{fontenc}
\usepackage{lmodern}

% Imágenes
\usepackage[pdftex]{graphicx}
\usepackage{latexsym}
\usepackage{fancybox}

% Ruta para las imágenes
\graphicspath{{img/}}

% Rotaciones
\usepackage[twoside]{rotating}

% Multirow
\usepackage{multirow}

% Referencias
\usepackage[spanish]{varioref}
\usepackage[pdftex,colorlinks=true,linkcolor=black]{hyperref}

% Colores
\usepackage{color}
\usepackage{colortbl}
\definecolor{rojo}{rgb}{1.0,0.0,0.0}

% Párrafos
\setlength{\parskip}{6pt}

% Code for creating empty pages
% No headers on empty pages before new chapter
\makeatletter
\def\cleardoublepage{\clearpage\if@twoside \ifodd\c@page\else
    \hbox{}
    \thispagestyle{empty}
    \newpage
    \if@twocolumn\hbox{}\newpage\fi\fi\fi}
\makeatother \clearpage{\pagestyle{empty}\cleardoublepage}

%% Entorno Listings para código fuente
\usepackage{listingsutf8}[2007/11/11]

\lstset{
  frame=Ltb, %forma del cuadro
  framerule=0pt, %ancho del borde
  aboveskip=0.5cm, %separación de los números de línea
  framexleftmargin=0.4cm, %margen externo izquierdo
  framesep=0pt,
  rulesep=.4pt,
  rulesepcolor=\color{black},
  % 
  stringstyle=\ttfamily,
  showstringspaces = false,
  basicstyle=\scriptsize,
  keywordstyle=\bfseries,
  % 
  numbers=left,
  numbersep=15pt,
  numberstyle=\tiny,
  numberfirstline= false,
  %
  inputencoding=utf8/latin1
}

% minimizar fragmentado de listados
\lstnewenvironment{listing}[1][]{
  \lstset{#1}\pagebreak[0]}{\pagebreak[0]
}


% Fancyhdr - Encabezados y pies de página
\usepackage{fancyhdr}
% Márgenes
\headsep=8mm
\footskip=14mm

% Fancy Header Style Options
\pagestyle{fancy}               % Sets fancy header and footer
\fancyfoot{}                    % Delete current footer settings

% Sin mayúsculas en la cabecera
\lhead{\nouppercase{\rightmark}}
\rhead{\nouppercase{\leftmark}}

% Capítulo
\renewcommand{\chaptermark}[1]{ % Lower Case Chapter marker style
  \markboth{\chaptername\ \thechapter.\ #1}{}} 

% Sección
\renewcommand{\sectionmark}[1]{ % Lower case Section marker style
  \markright{\thesection.\ #1}} 

% Página
\fancyhead[LE,RO]{\bfseries\thepage} % Page number (boldface) in left on even
                                     % pages and right on odd pages

% ------------------ Macro para encabezados y pies de página-------------------
%    Uso: \encabezado{pares(pag izquierda)}
% -----------------------------------------------------------------------------
\def\encabezado{
  \fancyhead[RE]{\bfseries\leftmark}     % In the right on even pages
  \fancyhead[LO]{\bfseries\rightmark}  % In the left on odd pages
  \renewcommand{\headrulewidth}{0.5pt} % Width of head rule
}
% -----------------------------------------------------------------------------

% ------------------ Macro para insertar una imagen ---------------------------
%    Uso: \imagen{nombreFichero}{Ancho(cm))}{Etiqueta}{Identificador}
% -----------------------------------------------------------------------------
\usepackage{float}
\usepackage{ifthen}
\def\imagen#1#2#3#4{
  \begin{figure}[H]
    \begin{center}
      \ifthenelse{\equal{#2}{}}
      {\includegraphics{#1}}{\resizebox{#2cm}{!}{\includegraphics{#1}}}
      \ifthenelse{\equal{#3}{}}{}{\caption{#3}}
      \label{#4}
    \end{center}
  \end{figure}
}
% -----------------------------------------------------------------------------


% ------------------ Macro para la portada ------------------------------------
%    Uso: \portada{asignatura}{titulo}{subtítulo}{autor}{fecha}
% -----------------------------------------------------------------------------
\def\portada#1#2#3#4#5{
  \thispagestyle{empty}
  \vspace*{-3.3cm}

  \begin{minipage}[t]{14cm}
    \begin{center}
      \includegraphics[scale=0.25]{logoesi}
  
      \vspace*{1.5cm}
      {\Large \textbf{Universidad de Castilla-La Mancha\\ 
          Escuela Superior de Informática}\\}
    
      \vspace*{1.2cm}
      {\huge \textbf{#1}\\}
    
      \vspace*{1.5cm}
      {\huge #2}\\{\Large #3\\}
    
      \vspace*{1.5cm}
      {\large #4\\}
      \vspace*{1.4cm}
      \large{#5}
    \end{center}
  \end{minipage}

  \newpage
%   \vspace*{1cm}
%   \thispagestyle{empty} 
%   \newpage
}
% -----------------------------------------------------------------------------


% ------------------ Macro para la licencia -----------------------------------
%    Uso: \portada{asignatura}{titulo}{subtítulo}{autor}{fecha}
% -----------------------------------------------------------------------------
\def\licencia#1{
  \thispagestyle{empty}  % Suprime la numeración de esta página
  \vspace*{16.5cm}
  \begin{small}
    \copyright~ #1. Se permite la copia, distribución y/o 
    modificación de este documento bajo los términos de la licencia de
    documentación libre GNU, versión 1.1 o cualquier versión posterior publicada
    por la {\em Free Software Foundation}, sin secciones invariantes. Puede
    consultar esta licencia en http://www.gnu.org. \\[0.2cm]
    Este documento fue compuesto con \LaTeX{}. 
  \end{small}
  \newpage
%   \thispagestyle{empty}
%   \vspace*{0cm}
%   \newpage
}
% -----------------------------------------------------------------------------




% Encabezado y pie de página
\encabezado

\begin{document}

% Silabación extra
\hyphenation{
a-sig-na-tu-ras
au-to-ma-ti-za-rá
ca-tá-lo-go
ca-rre-ra
diag-nos-tico
in-fe-rior
man-te-ni-mien-to
per-so-nal
pro-por-cio-na-rá
pu-bli-ca-da
re-qui-si-tos
res-pecto
u-su-a-rios
vi-lla-rre-al
}


% Portada
\portada{Planificación y Gestión de\\Sistemas de Información}
{Trabajo 5}{Estimación del software}
{Sergio de la Rubia García-Carpintero\\Miguel Millán Sánchez-Grande\\
  Luis Muñoz Villarreal\\Alicia Serrano Sánchez\\
  Juan Miguel Torres Triviño}{24 de Mayo de 2010}

% Licencia
\licencia{Sergio de la Rubia García-Carpintero, Miguel Millán Sánchez-Grande,
  Luis Muñoz Villarreal, Alicia Serrano Sánchez, Juan Miguel Torres Triviño}

\chapter*{Ficha de trabajo}
\begin{description}
\item[Código] T5
\item[Fecha] 24 de Mayo de 2010
\item[Título] Estimación del software
\end{description}

\begin{table}[!ht]
  \centering
  \begin{tabular}{lp{5cm}c}
    \multicolumn{3}{l}{\Large \textbf{Equipo} G4} \\ \\
    \multicolumn{1}{c}{\emph{Apellidos y nombre}} & 
    \multicolumn{1}{c}{\emph{Firma}} & \emph{Puntos} \\
    \hline \\
    de la Rubia García-Carpintero, Sergio & & 8 \\ \\
    Millán Sánchez-Grande, Miguel         & & 8 \\ \\
    Muñoz Villarreal, Luis                & & 8 \\ \\
    Serrano Sánchez, Alicia               & & 8 \\ \\
    Torres Triviño, Juan Miguel           & & 8 \\ \\
    \hline
  \end{tabular}
\end{table}

% Índices
\tableofcontents
\listoftables
%\listoffigures

%% INICIO DEL DOCUMENTO %%%%%%%%%%%%%%%%%%%%%%%%%%%%%%%%%%%%%%%%%%%%%%%%%
\chapter*{Introducción}
%breve introducción del software llamado...

\section{Visión general}

\section{Objetivos}

\chapter{Identificación de los módulos}
\section{Gestión de usuarios}
\subsection{Entradas externas}
\subsubsection{Identificación del usuario}
Cada usuario debe identificarse mediante un nombre y una contraseña.

\subsubsection{Preferencia de horarios} 
Cada profesor introduce sus preferencias de horarios de trabajo para que el
sistema lo tenga en cuenta a la hora de la realización del horario oficial
del curso.
\subsection{Salidas externas}
\subsubsection{Datos personales}
Muestra los datos personales de un determinado profesor, incluyendo
nombre, correo electrónico, teléfono, publicaciones, currículum, \dots
\subsubsection{Horario de un profesor}
El sistema muestra el horario que tiene asignado un determinado profesor.

\subsubsection{Tutorías de los profesores} 
Muestra los horarios de tutorías junto con el despacho asignado de los
profesores de cada facultad.
 
\subsubsection{Asignaturas impartidas por cada profesor}
El sistema muestra la información de cada asignatura que comprende:
\begin{itemize}
\item Profesor/es que la imparte/n.
\item Créditos.
\item Contenidos de la asignatura.
\item Asignaturas recomendadas.
\item Página web de la asignatura.
\item Planificación docente.
\item Sistema de evaluación.
\end{itemize}

\subsection{Consultas externas}
\subsubsection{Consultar currículum}

Información del currículum de un determinado profesor.

\subsubsection{Consultar datos de asignatura} 
Información sobre una determinada asignatura (créditos, profesor que la
imparte, planificación, \dots
\subsubsection{Horarios de asignaturas}
Horarios de las asignaturas de cada curso. Cada horario
también muestra información de la localización en la que cada asignatura se
imparte. 

\subsection{Archivos lógicos de interfaz externos}
\subsubsection{Datos de los profesores}
El sistema tiene una base de datos donde se guarda toda la información con
respecto a los profesores (nombre, apellidos, teléfono, correo electrónico,
asignaturas que imparte, \dots)

\subsubsection{Datos de las asignaturas}
En la base de datos de la universidad se tienen los datos de las asignaturas,
tales como nombre, créditos, guión de la asignatura, contenidos, evaluación,
\dots 


\section{Gestión de cursos}
\subsection{Entradas externas}
\subsubsection{Añadir curso}
Entrada destinada a la inclusión de los distintos cursos correspondientes al
plan de estudios en el sistema.

\subsubsection{Añadir asignaturas}
Entrada generada para la inserción de las diferentes asignaturas asociadas a
cada curso.

\subsubsection{Añadir aulas}
Entrada destinada a la introducción las aulas disponibles destinadas al
correspondiente plan de estudio.

\subsubsection{Añadir profesores}
Entrada generada para la incluir los profesores encargados de la docencia de
cada asignatura.

\subsection{Salidas externas}
\subsubsection{Relación asignaturas/profesores}
Información 
\subsubsection{Relación cursos/aulas}
\subsubsection{Cuadrante de horarios}

\subsection{Consultas externas}
\subsubsection{Consultar relación asignaturas/profesores}
\subsubsection{Consultar relación cursos/aulas}
\subsubsection{Consultar cuadrante de horarios}

\subsection{Archivos lógicos internos}
\subsubsection{Datos de las asignaturas}
\subsubsection{Datos de las cursos}
\subsubsection{Datos de las aulas}

\subsection{Archivos lógicos de interfaz externos}
\subsubsection{Historial de guías docentes}

\section{Interfaz}
\subsection{Entradas externas}
\subsubsection{Autentificación de los usuarios}
Recoge los datos necesarios para ver si un usuario entra al sistema,
dependiendo si éste es administrador o no.

\subsubsection{Introducción de datos para la inserción/modificación/borrado 
de información de un determinado usuario}
Por su similitud se explican juntas pero equivale a tres \emph{funciones de
usuario} diferentes. El objetivo de estas funciones es recoger los datos
necesarios para la de inserción, modificación o borrado de la información de
un usuario en la aplicación.

\subsubsection{Introducción de datos para la inserción/modificación/borrado 
de información de un determinado curso }
Por su similitud se explican juntas pero equivale a tres \emph{funciones de
usuario} diferentes. El objetivo de estas funciones es recoger los datos
necesarios para la de inserción, modificación o borrado de la información de
un curso en la aplicación.

\subsection{Salidas externas}
\subsubsection{Mostar los datos para una búsqueda de información con
posibilidad de ordenación por diferentes criterios}
El objetivo de este proceso es mostrar un listado de usuarios que cumplen las
condiciones de búsqueda solicitadas y en el orden indicado.

\subsubsection{Mostrar los datos de un determinado usuario}
El objetivo de este proceso es mostrar información detallada referente a un
determinado usuario (profesores).

\subsubsection{Mostrar los datos de un determinado curso}
El objetivo de este proceso es mostrar información detallada referente a un
curso, incluyendo horario, aulas, asignaturas, etc.

\subsection{Consultas externas}
\subsubsection{Búsqueda de información con posibilidad de ordenación por
  diferentes criterios}
El objetivo de este proceso es la introdución de unos criterios para realizar 
una búsqueda detallada tanto de cursos como de usuarios.

\subsubsection{Consulta de los datos de un determinado usuario}
El objetivo de esta función es la consulta simple de la información referente
a un usuario (profesor).

\subsubsection{Consulta de los datos de un determinado curso}
El objetivo de esta función es la consulta simple de la información referente
a un curso.

\subsection{Archivos lógicos internos}
\subsubsection{Importación de datos desde un fichero}
El objetivo de esta función es introducir los datos necesarios para importar 
inserciones, modificaciones o borrados desde un fichero de texto.

\subsubsection{Exportación de los resultados de la búsqueda a un fichero}
El objetivo de esta función es introducir los datos para obtener de la base
de datos de la aplicación el listado de usuarios que cumplen las condiciones de
búsqueda solicitadas y en el orden indicado en un fichero de texto.

\subsection{Archivos lógicos de interfaz externos}
No se ha detectado ningún archivo lógico de interfaz externo.

\chapter{Determinación de la complejidad de las funciones de los 
  módulos}
\section{Gestión de usuarios}
\subsection{Entradas externas}
\subsubsection{Identificación del usuario}
La complejidad de esta función de usuario es \textbf{baja} porque se hace
referencia a un registro para la información de usuario y el número de datos
elementales es dos. 

\subsubsection{Preferencia de horarios} 
La complejidad de esta función es \textbf{media} porque se hace
referencia a dos registros donde se tienen guardados los horarios y las
asignaturas y el número de datos elementales son que se utilizan son seis. 

\subsection{Salidas externas}
\subsubsection{Datos personales}
La complejidad de esta función es \textbf{media} porque se hace referencia a
un registro para la información de usuario y los datos elementales que se
utilizan son dieciséis.

\subsubsection{Horario de un profesor}
La complejidad de esta función es \textbf{media} porque se hace referencia a
dos registros, horarios y asignaturas; y los datos elementales que
se utilizan son seis.

\subsubsection{Tutorías de los profesores} 
La complejidad de esta función es \textbf{baja} porque se hace referencia a
dos registros, profesores y tutorías; y el número de datos elementales que se
utilizan son cuatro. 

\subsubsection{Asignaturas impartidas por cada profesor}
La complejidad de esta función es \textbf{medio} porque se hace referencia a
dos registros, profesores y asignaturas; y el número de datos elementales que se
utilizan son siete. 

\subsection{Consultas externas}
\subsubsection{Consultar currículum}
La complejidad de esta función es \textbf{baja} porque se hace referencia a
un registro y el número de datos elementales es dos.

\subsubsection{Consultar datos de asignatura} 
La complejidad de esta función es \textbf{baja} porque se hace referencia a
un registro y el número de datos elementales es nueve.

\subsubsection{Horarios de asignaturas}
La complejidad de esta función es \textbf{baja} porque se hace referencia a
tres registros y los datos elementales que se utilizan son cuatro.

\subsection{Archivos lógicos de interfaz externos}
\subsubsection{Datos de los profesores}
La complejidad de esta función es \textbf{baja} porque se hace referencia a
un registro y los datos elementales son siete.

\subsubsection{Datos de las asignaturas}
La complejidad de esta función es \textbf{baja} porque se hace referencia a
un registro y los datos elementales son catorce.

\section{Gestión de cursos}
\subsection{Entradas externas}
\subsection{Salidas externas}
\subsection{Consultas externas}
\subsection{Archivos lógicos internos}
\subsection{Archivos lógicos de interfaz externos}

\section{Interfaz}
\subsection{Entradas externas}
\subsubsection{Autentificación de los usuarios}
La complejidad de esta función de usuario es \textbf{baja}, porque tiene un
registro para la información del usuario y el número de tipo de datos
elementales es dos. 

\subsubsection{Introducción de datos para la inserción/modificación/borrado 
de información de un determinado usuario}
La complejidad para cada función de usuario (inserción, modificación y
borrado) es \textbf{media}, porque tiene un registro para la información del
usuario y el número de tipo de datos elementales es diecisiete.

\subsubsection{Introducción de datos para la inserción/modificación/borrado 
de información de un determinado curso }
La complejidad para cada función de usuario (inserción, modificación y
borrado) es \textbf{baja}, porque tiene un registro para la información del
curso y el número de tipo de datos elementales es diez.

\subsection{Salidas externas}
\subsubsection{Mostar los datos para una búsqueda de información con
posibilidad de ordenación por diferentes criterios}
La complejidad de esta función de usuario es \textbf{media}, porque tiene un
registro para la información y el número de tipo de datos elementales es veinte.

\subsubsection{Mostrar los datos de un determinado usuario}
La complejidad de esta función de usuario es \textbf{baja}, porque tiene un
registro para la información del usuario y el número de tipo de datos
elementales es nueve.

\subsubsection{Mostrar los datos de un determinado curso}
La complejidad de esta función de usuario es \textbf{baja}, porque tiene un
registro para la información del curso y el número de tipo de datos
elementales es cinco.

\subsection{Consultas externas}
\subsubsection{Búsqueda de información con posibilidad de ordenación por
  diferentes criterios}
La complejidad de esta función de usuario es \textbf{media}, porque no tiene
ningún registro para la información y el número de tipo de datos elementales
es dieciocho.

\subsubsection{Consulta de los datos de un determinado usuario}
La complejidad de esta función de usuario es \textbf{baja}, porque no tiene
ningún para la información del usuario y el número de tipo de datos
elementales es uno.

\subsubsection{Consulta de los datos de un determinado curso}
La complejidad de esta función de usuario es \textbf{baja}, porque no tiene
ningún registro para la información del curso y el número de tipo de datos
elementales es uno.

\subsection{Archivos lógicos internos}
\subsubsection{Importación de datos desde un fichero}
La complejidad de esta función de usuario es \textbf{baja}, porque tiene un
fichero referenciado para la información y el número de tipo de datos
elementales es nueve.

\subsubsection{Exportación de los resultados de la búsqueda a un fichero}
La complejidad de esta función de usuario es \textbf{alta}, porque tiene dos
fichero referenciado para la información y el número de tipo de datos
elementales es veinte.

\chapter{Tablas de puntos función sin ajustar}
\section{Interfaz}
\begin{table}[!h]
  \centering
  \begin{tabular}{|b{2.1cm}<\centering|b{2.1cm}<{\centering}|c|c|c|c|c|}
    \hline
    \textbf{Tipo de función} & \textbf{Nivel de complejidad} &
    \textbf{Número} & \textbf{*} & \textbf{Peso} & \textbf{=} & \textbf{Total}\\
    \hline \hline
    Entradas  &  &  &  &  &  &  \\ \hline
    Salidas &  &  &  &  &  &  \\ \hline 
    Consultas &  &  &  &  &  &   \\ \hline
    Archivos &  &  &  &  &  &   \\ \hline
    Interfaces &  &  &  &   &  &\\ \hline
  \end{tabular}
  \caption{Puntos de función sin ajustar del módulo interfaz} 
  \label{Tab:PFSAint}
\end{table}
\chapter{Explicación de los valores de los factores de influencia}
\chapter{Cálculo de los puntos función ajustados}

\end{document}
